\section{Natural compatibility of the three quotients}\label{sec:unification}
The preceding binary theorem determines a metric from a pre-existing
observation mechanism. The next statement identifies the associated
root sheets and endpoint profile without an embedding into a designed
contrast model. The distinction between labelled roots and their
symmetric observation is important here.

\begin{theorem}[Calibrated binary structural compatibility]\label{thm:unification}
Fix a compact separated configuration in
Theorem~\ref{thm:binary-hankel}, with positive cell and exposure margins.
Near its repeated-root centre, the polynomial map from labelled roots to
cluster sums, variances and inward endpoint shifts is finite. Over
strictly positive variances its degree is $2^k$, corresponding exactly
to exchange of the two roots in each interior cluster.

For a fixed retained direction $w=(x,z)\in\R_+^{k+e}$, the complete leading
observation residual at scale $t$, with unrestricted free sum and mixing
directions $a\in\R^{k+1}$, is the affine sheet
\begin{equation}\label{eq:binary-leading-sheet}
 E_w=\{\mathsf La+\mathsf Dw:a\in\R^{k+1}\}.
\end{equation}
Here the probability score blocks are obtained from \eqref{eq:binary-score}
by multiplying row $j$ by the four-vector
$\pi_j(1-\pi_j)(M_1-M_2)$. The root labels give identical observation
sheets, so their multiplicity is removed in the intrinsic residual set.
Its metric tensor envelope is the affine-sheet envelope of
Theorem~\ref{thm:lift} whenever $w\ne0$; for $w=0$ the minimum is zero.

The following three operations therefore give compatible objects derived
from the same model: the finite root quotient and its leading residual,
profiling the free coordinates to obtain $Q$ with the native Hankel
inverse \eqref{eq:hankel-cone}, and profiling the inward endpoints to
obtain $G_Q(x)$. More precisely, locally at the fixed centre,
\begin{equation}\label{eq:binary-variational}
 \lim_{t\downarrow0}\frac4{t^2}
 \inf_{\text{free sums, mixing, endpoints}}
 \sum_j\lambda_j h^2(p_j(t),p_j(0))=G_Q(x),
\end{equation}
where the interior variances are fixed to $t x$, and the other variables
range over their feasible local neighbourhood. The endpoint-free version
with all of $w$ fixed equals $w^{\mathsf T}Qw$.
These limits and the score expansion are uniform for bounded retained
directions and on compact separated centre sets. In the one-colour
subclass, the endpoint dimensions and fibres are exactly those in
Theorem~\ref{thm:native-visibility}.
\end{theorem}
\begin{proof}
For a labelled interior pair with displacements $u_1,u_2$, write
$s=u_1+u_2$ and $v=(u_1-u_2)^2/4$. Its factor is exactly
$(T-r-s/2)^2-v$. The inverse is
$u_1=s/2+\sqrt v$, $u_2=s/2-\sqrt v$, with exchange as the only other
label. Distinct separated clusters cannot mix in a sufficiently small
neighbourhood. Endpoint coordinates are their inward displacements.
This proves finiteness, the stated degree, and feasibility at $v=0$ by
continuity; all small nonnegative variances are realized.

In the symmetric coordinates, the map to the vector of clock log odds
is analytic and has derivative $V$, which is invertible by
Theorem~\ref{thm:binary-hankel}. It is consequently a local analytic
diffeomorphism on an open extension of those coordinates. Each calibrated
probability table determines its scalar mixing probability by any cell
where $M_1-M_2$ is nonzero. Strict probability margins make these
coordinate changes and their derivatives uniformly bounded. The
observation map is therefore locally bi-Lipschitz in the symmetric
coordinates. This statement concerns the chosen local germ, not a
uniform inverse through changes of cluster pattern or channel rank.

Set the symmetric coordinates to $(t a,t w)$. Taylor expansion gives
$p(t)=p(0)+t(\mathsf La+\mathsf Dw)+O(t^2)$ uniformly on bounded sets.
Conversely, a bounded observation residual divided by $t$ forces all
symmetric coordinate increments to be $O(t)$ by the local inverse.
Taking subsequential limits, and realizing arbitrary bounded $a$, proves
both inclusions in \eqref{eq:binary-leading-sheet}. The full score is
injective, so for nonzero $w$ the sheet does not contain zero and its
augmented affine matrix has full rank, as required in
Theorem~\ref{thm:lift}.

For the minimization in \eqref{eq:binary-variational}, zero free increments
and zero endpoint increments give a trial of cost $O(t^2)$ for bounded
$x$. Norm equivalence between Hellinger distance and probability
increments, followed by the local inverse, forces every minimizing
sequence to have all symmetric coordinate increments $O(t)$. Positive
definiteness of the score Gram matrix also bounds the scaled minimizers
uniformly. On these bounded sets,
\[
 \frac4{t^2}\sum_j\lambda_j h^2(p_j(t),p_j(0))
 =\|\mathsf La+\mathsf D(x,z)\|_W^2+o(1),
 \qquad W=\diag(\lambda_j/p_{jc}).
\]
The minimum over $a$ is the nuisance Schur complement, and the remaining
minimum over $z\ge0$ is exactly $G_Q(x)$. The unrestricted free minimizer
and the nonnegative endpoint minimizer are feasible in the local chart
for small time. Thus testing them gives the upper bound; localization
and the uniform expansion give the lower bound. The same argument with
$w$ fixed gives its quadratic cost. Compact separation and positivity
bounds supply all the asserted uniform constants.
\end{proof}
This compatibility theorem does not claim that distinct labelled root
sheets must yield distinct residual sheets. In this natural symmetric
model they do not. Nor does calibration of the channels identify the
unknown-channel experiment with the calibrated one. The general
finite-map results, the paired benchmark, and the preserved
unknown-channel binary inverse retain their separate hypotheses.

\subsection*{Preparation and verification}
The development and drafting of the residual-scale and divisorial
criteria in Section~\ref{sec:residual-scale}, the one-cell and two-endpoint
fibre arguments in Section~\ref{sec:walls}, and the interpolation-based
Hankel classification in Sections~\ref{sec:binary-hankel}--\ref{sec:unification}
used language-model assistance. Their proofs are supplied in the article.
The accompanying exact algebra checks test specified identities and finite
instances; they do not certify the universal theorems. This revision is
prepared for independent mathematical and author review.
